\chapter{Institutions for Long Lives}
\label{ch:30}

A society whose members lived for centuries would differ from ours in the structure of its institutions more than in the details of its technology, and the differences would not be anticipated by extrapolating from present practice. Nearly every institution that humans have devised presupposes a lifespan. Property passes by inheritance because owners die; offices rotate because officeholders retire or die; marriage, citizenship, apprenticeship, and the criminal law all assume stages of life that last decades. The assumption of mortality does silent work in preventing the accumulation of advantage, in refreshing the ranks of those who decide, and in allowing the world to be inherited by persons who did not choose it. Max Planck is reputed to have said that science advances one funeral at a time, and whether or not he said it, the remark identifies a mechanism by which societies renew themselves that a long-lived population would remove. The thesis of this chapter is that extended life, if it were achieved, would require deliberate substitutes for these mechanisms, that the substitutes can be designed, and that their design is an institutional problem whose features can be stated now with some precision.

Begin with demography, because it fixes the scale of the problem. In a population of constant size, births must balance deaths, and if the mortality rate is the reciprocal of a mean lifespan $L$, the population is the product of the annual number of births and the mean lifespan. A world of eight billion persons with a mean lifespan of eighty years has some hundred million births a year, as is roughly the case. If the mean lifespan were a thousand years and the population were to remain the same, births would have to fall to eight million a year, one-twelfth of the present number; at ten thousand years, to eight hundred thousand; and at a million years, to eight thousand. A population that wished to grow while living longer would have to accept still fewer births relative to its size, or would have to expand into new habitats, and a population that did not control births would grow at a rate that no environment could sustain. The consequences touch the most intimate decisions of family and the most contested questions of the right to reproduce. A program that extended life without addressing them would be proposing to resolve a constraint by default, and the default would be either an explosion of numbers or a rationing of births by whoever controlled access to the means of extension. It is part of the argument of the book that neither is acceptable, and that the rate of births is a matter for the persons affected through procedures they can contest.

The age structure of a long-lived population has implications for representation. In a stationary population with a constant mortality rate, the proportion of persons older than $a$ is $e^{-a/L}$. For a mean lifespan of a thousand years, $61$ per cent of the population would be older than five hundred years and only $7.7$ per cent younger than eighty. Most of the adults would be old in a sense that has no present counterpart, and in a society that gave votes and offices on the basis of age or of years of residence the young would be a permanent minority. The arrangements for voting, officeholding, and the allocation of public resources would have to be designed with this in view, and the principle that should govern them is that the participation of persons in decisions about their future should not depend on the accident of when they were born. This favors mechanisms that give weight to the interests of the young and to those not yet born, such as institutions with a mandate to represent them, and it counsels against rules that confer authority by age.

Wealth and power are the second concern. Capital, left to accumulate, grows at the rate of return, and an owner who does not die keeps all of the growth. If the rate of return exceeds the growth rate of the economy by $r-g$, the share of the total held by a perpetual owner of fixed initial wealth increases in proportion to $e^{(r-g)T}$. At a difference of two per cent, the multiple is $7.4$ after a century, $22{,}000$ after five hundred years, and nearly five hundred million after a thousand. These figures reflect the mathematics of compound growth, and not a forecast: no actual fortune has grown without interruption by taxation, expropriation, war, and misjudgment, and the fortunes of dynasties have regularly dissipated. But the dissipation has been largely a consequence of mortality, which subjects wealth to inheritance, division, and the unequal abilities of heirs. A population of immortals who made good investments would not be subject to these processes, and the equilibrium of wealth in such a society would be one of extreme concentration unless institutions prevented it. The instruments available are familiar. A periodic levy on wealth at a rate not less than $r-g$ prevents the share of any holder from growing without bound; limits on the duration of property rights, as in the law of patents and copyright, limit the accumulation of rent; and a requirement that wealth be returned to the common stock in part at stated intervals replicates for the living what inheritance taxation does for the dead. Equally important, the same reasoning applies to political power, which accumulates by the same mechanism of the compounding of advantage, and the terms of offices, the rotation of those who hold them, and the opening of positions to the young are the political counterpart of the levy on wealth.

The third concern is the durability of commitments. A government that binds itself for centuries is exerting a dead hand over those who come after it, and the dead-hand problem, familiar from the law of trusts and the debate over constitutional entrenchment, would be altered in a long-lived society, because the persons who made the commitments would be among those who lived under them. This has some advantages: the makers of a rule can be held to account for its consequences, and would have an incentive to revise a rule that did not work. It also has dangers, since commitments made by powerful persons could be sustained against challenge by the same persons for as long as they lived. The analysis of Chapter 28 suggests a principle for the interval of renewal. A person's consent, and a body's mandate, decay with the connectedness of the stages that gave it, and the interval over which they can be presumed to hold is of the order of the window of strong connectedness, some decades. Commitments of greater duration should lapse unless renewed, with the renewal carried out by the persons then affected, who include persons not alive at the time of the original decision. A sunset provision of this kind is the mildest form of the principle, and it is the one most readily combined with the existing practices of legislatures.

The fourth concern, and the most difficult, is access. A technology of extended life would not at first be available to all, and the early distribution would be determined by cost, geography, and political influence. The history of medical advances is that those which are scarce are rationed by wealth, and that the gap between rich and poor in life expectancy, already large, would be widened. The program's principle of the admissible range requires that the distribution of such a technology be a matter for deliberate decision by procedures affected persons can contest, and the criteria of a just distribution are not for this book to determine. The book can say what is not acceptable: a distribution settled by the ability to pay without any public procedure, and a technology whose benefits are confined to the persons who are most able to influence the decisions about it. It can also say what must be provided in any scheme: a public record of who has access and on what terms, an independent audit of the effects on mortality and health, and the guarantee, already stated in Chapter 28, that the declining of the technology does not result in disadvantage.

The final concern is the continuing authority of the persons who hold the program together. A civilization that lasts for millions of years will need institutions that last as long, and the history of institutions is not encouraging: the longest-lived organizations in the human record, certain religious orders, a few commercial enterprises and ruling houses, have survived for centuries and, in a handful of cases, for more than a thousand years. Their persistence has depended on features that can be identified: a clear and limited purpose, a method for the training and selection of successors, a record that can be consulted, a capacity to change the rules while remaining the same body, and a culture of accountability. These are the features that the preceding chapters require of the institutions of verification and authority. What the programs of the book add is the requirement that the institutions be maintained for periods in which the survival of a single one of them is not to be assumed, which points again to redundancy, to diversity of design, and to the capacity to rebuild an institution from its records, and not to the hope that a single institution will persist unchanged.

\section*{Formalization}

Three results are given, for population, for accumulation, and for renewal.

Consider a population of constant size $N$ with birth rate $B$ per year and a constant mortality hazard $\mu=1/L$ at all ages (the extrinsic-hazard limit of Chapter 27), so that survival to age $a$ is $e^{-\mu a}$.

\begin{proposition}
The stationary population satisfies $N=B/\mu=BL$. The age distribution is exponential with density $\mu e^{-\mu a}$, so the proportion over age $a$ is $e^{-a/L}$ and the proportion below age $a$ is $1-e^{-a/L}$. To maintain a fixed $N$ as $L$ rises the birth rate must fall as $B=N/L$.
\end{proposition}

\begin{proof}
The balance $dN/dt=B-\mu N=0$ gives $N=B/\mu$. The age density is proportional to the survival function $e^{-\mu a}$ and is normalized by $1/\int_0^\infty e^{-\mu a}da=\mu$.
\end{proof}

For $N=8\times10^9$: $B=10^8$, $8\times10^6$, $8\times10^5$, $8\times10^3$ per year for $L=80$, $10^3$, $10^4$, $10^6$ years. The proportion younger than $80$ years is $0.632$ at $L=80$, $0.077$ at $L=10^3$, $0.0080$ at $L=10^4$, and $8\times10^{-5}$ at $L=10^6$. The proportion older than $500$ years is $0.607$ at $L=10^3$ and $0.951$ at $L=10^4$. (The exponential law is the extrinsic-hazard limit and is used because it is the case in which senescence is removed; under it the mean age of the living is $L$.)

For accumulation, let an owner hold wealth $W(t)$ growing at rate of return $r$ net of consumption, in an economy of total wealth growing at rate $g$, with a levy at rate $\lambda$ per year on the owner's wealth.

\begin{proposition}
The owner's share of total wealth is $s(T)=s(0)e^{(r-g-\lambda)T}$ while $s\ll1$. The share is bounded for all $T$ if and only if $\lambda\ge r-g$. Without a levy, the time for the share to multiply by a factor $K$ is $T_K=\ln K/(r-g)$.
\end{proposition}

\begin{proof}
$\dot W=(r-\lambda)W$ and the total grows as $e^{gt}$, so the ratio grows at $r-\lambda-g$. Boundedness requires a nonpositive exponent. The doubling-type time follows by solving $e^{(r-g)T_K}=K$.
\end{proof}

For $r-g=0.02$ the share multiplies by $K=7.39$, $2.2\times10^4$, $4.85\times10^8$ in $T=100$, $500$, $1000$ years; it doubles in $34.7$ years; and a levy of $2$ per cent holds the share constant. The result holds for the share while it remains small, and the qualification is relevant to the claim of extreme concentration: the model becomes inaccurate as $s\to1$, at which point the owner's own return affects $r$.

For renewal, suppose a commitment made at time $0$ by a body whose stages decay in connectedness as in Chapter 28, with mandate strength $c(t)=e^{-t/\tau}$. A commitment is renewed whenever $c(t)$ falls to the threshold $\theta$.

\begin{proposition}
The renewal interval is $\Delta_R=\tau\ln(1/\theta)$, and over a horizon $T$ the number of renewals required is $\lfloor T/\Delta_R\rfloor$. For $\tau=20$ years and $\theta=1/2$, $\Delta_R=13.9$ years, and a commitment intended to last $10^3$ years requires $72$ renewals.
\end{proposition}

\begin{proof}
Solve $e^{-\Delta_R/\tau}=\theta$.
\end{proof}

The requirement is unusually burdensome compared with the practice of legislation, which renews few things at such intervals, and it is intended as a limit case; a more tolerant threshold $\theta=0.1$ gives $\Delta_R=46$ years and $21$ renewals for a thousand years. The choice of $\theta$ is a normative one, to be made by the society concerned, and the calculation shows how the burden of renewal depends on it.
